% Generated draft: requires mathematical and source review.
\documentclass[11pt]{article}
\usepackage[margin=25mm]{geometry}
\usepackage{fontspec}
\setmainfont{Latin Modern Roman}
\usepackage{amsmath,amssymb,mathtools,unicode-math}
\setmathfont{Latin Modern Math}
\usepackage{enumitem,needspace,xurl,hyperref}
\setlength{\parindent}{0pt}
\setlength{\parskip}{5pt}
\setlength{\emergencystretch}{3em}
\newcommand{\N}{\mathbb{N}}
\newcommand{\Z}{\mathbb{Z}}
\newcommand{\Q}{\mathbb{Q}}
\newcommand{\R}{\mathbb{R}}
\newcommand{\C}{\mathbb{C}}
\newcommand{\F}{\mathbb{F}}
\newcommand{\A}{\mathbb{A}}
\newcommand{\PP}{\mathbb{P}}
\newcommand{\E}{\mathbb{E}}
\newcommand{\eps}{\varepsilon}
\newcommand{\dd}{\,\mathrm{d}}
\begin{document}
\section*{UnsolvedMath ID: 800639}
\textbf{Source problem code:} catalogue entry
\subsection{Definitions and mathematical statement}
Let \(\N=\{0,1,2,\ldots\}\). For \(f,g\in\N^{\N}\), write \(f\leq^{*}g\) if \(f(n)\leq g(n)\) for all but finitely many \(n\). The bounding and dominating numbers are
\[
\mathfrak b=\min\bigl\{|B|:B\subseteq\N^{\N},\ (\forall g\in\N^{\N})(\exists f\in B)\ f\not\leq^{*}g\bigr\},
\]
\[
\mathfrak d=\min\bigl\{|D|:D\subseteq\N^{\N},\ (\forall f\in\N^{\N})(\exists g\in D)\ f\leq^{*}g\bigr\}.
\]
For an infinite cardinal \(\kappa\), \(\operatorname{cf}(\kappa)\) denotes its cofinality.

The supplied excerpt does not specify the underlying class of spaces, the ambient space for unions, or the conventions for covers. The following is a conventional countable-open-cover expansion of its notation, conditional on those conventions being the intended ones. For a topological space \(X\), let \(\mathcal O(X)\) consist of countable families \(\mathcal U\) of open proper subsets of \(X\) such that \(\bigcup\mathcal U=X\). Let \(\Gamma(X)\subseteq\mathcal O(X)\) consist of the infinite covers \(\mathcal U\) satisfying
\[
(\forall x\in X)\quad \bigl|\{U\in\mathcal U:x\notin U\}\bigr|<\aleph_0.
\]
The property \(\mathsf S_1(\Gamma,\Gamma)\) means that for every sequence \((\mathcal U_n)_{n\in\N}\) with \(\mathcal U_n\in\Gamma(X)\), there are choices \(U_n\in\mathcal U_n\) such that the family \(\{U_n:n\in\N\}\) belongs to \(\Gamma(X)\). In particular, the resulting family must be infinite.

For the background assertion, \(\mathsf U_{\mathrm{fin}}(\Gamma,\mathcal O)\) means that for every sequence \((\mathcal U_n)_{n\in\N}\) of members of \(\Gamma(X)\), each having no finite subcover, there are finite subfamilies \(\mathcal F_n\subseteq\mathcal U_n\) such that
\[
\left\{\bigcup\mathcal F_n:n\in\N\right\}\in\mathcal O(X).
\]

For a property \(P\) of admissible subspaces of a common ambient topological space, its additivity is conventionally
\[
\mathsf{add}(P)=\min\left\{|\mathcal F|:(\forall X\in\mathcal F)\ X\text{ has }P,\quad \bigcup\mathcal F\text{ does not have }P\right\},
\]
where the members and their union carry the relevant subspace topologies, and the minimum is taken over the intended admissible families. This expression presupposes that a witnessing family exists. The supplied record does not identify the ambient space or admissible class, so it does not determine that part of this definition.

The question attributed to Tsaban, Problem 11 (2003), is whether
\[
\mathsf{add}\bigl(\mathsf S_1(\Gamma,\Gamma)\bigr)=\mathfrak b.
\]
The source excerpt also reports the background bounds
\[
\mathfrak b\leq\mathsf{add}\bigl(\mathsf U_{\mathrm{fin}}(\Gamma,\mathcal O)\bigr)\leq\operatorname{cf}(\mathfrak d),
\]
and reports that the additivity of a corresponding combinatorial notion of smallness equals \(\mathfrak b\), citing \([\mathrm{AddQuad}]\). That combinatorial notion and the bibliographic expansion of this label are not supplied. These are source-reported background assertions; neither their current literature status nor their hypotheses have been independently verified here. The registry labels the question open.
\subsection{Short English statement}
Tsaban's Problem 11 asks whether the least number of spaces satisfying \(\mathsf S_1(\Gamma,\Gamma)\) whose union fails that property is the bounding number \(\mathfrak b\). The supplied excerpt leaves the intended class of spaces and cover conventions unspecified, so the expanded formulation is conditional on the stated conventions. The registry's open label is supplied metadata, not an independently verified current status.
\subsection{Sources}
\begin{itemize}
\item Local UnsolvedMath registry, permanent ID \(800639\), exact source problem code \texttt{catalogue entry}, title ``Selection Principles Milestone Problems --- Problem 11.'' The record attributes the problem to Tsaban, with proposed year 2003, and labels it open. Its supplied update date is 2026-10-16.
\item Supplied source attribution: Tsaban, \emph{Selection Principles in Mathematics: A Milestone of Open Problems} (2003), Problem 11, \url{https://arxiv.org/abs/math/0312182}. The record reports access on 2026-10-16 and extraction by \texttt{arxiv\_latex}; these are supplied provenance details, not a claim of access or verification in preparing this draft.
\item The source excerpt cites \([\mathrm{AddQuad}]\) for the assertion about a corresponding combinatorial notion of smallness. No full bibliographic entry, definition of that notion, or separate link is supplied.
\end{itemize}
\end{document}
